\RequirePackage[T1]{fontenc}
\documentclass[10pt,letterpaper]{article}
% Smoothing ladder of the best evolved design: built by `make smoothing-ladder` from
% ladder/ (receipt, generated table and level macros, re-encoded renders).
\usepackage{tgtermes,tgheros}
\usepackage[margin=0.6in]{geometry}
\usepackage{amsmath,graphicx,booktabs,array,xcolor,url}
\usepackage{pgfplots}
\pgfplotsset{compat=1.18}
\usepackage[hidelinks]{hyperref}
\hypersetup{pdftitle={Smoothing ladder of the best evolved design}, pdfauthor={Research draft; author approval pending}, pdflang={en-US}}
\definecolor{ink}{HTML}{132D3B}
\definecolor{teal}{HTML}{007E87}
\definecolor{rust}{HTML}{A3442D}
\newcommand{\code}[1]{{\small\nolinkurl{#1}}}
\graphicspath{{ladder/}}
\setlength{\parindent}{0pt}
\setlength{\parskip}{4pt}
\pagestyle{plain}
% \LadderLevel{level}{pin deg}{iterations}{F}{F with B-rep}{max mm}{rms mm}{tip SF}{mass g}
\newcommand{\LadderLevel}[9]{%
  \clearpage
  {\sffamily\bfseries\color{ink}\Large Level L#1\quad
   \ifnum#3=0 as exported (no smoothing)\else #3 iterations, feature pins #2$^\circ$\fi}\par
  {\sffamily $F=#4$\quad(with B-rep kept: $#5$)\quad
   max displacement #6\,mm,\ RMS #7\,mm\quad tip SF #8\quad 3D mass #9\,g}\par\smallskip
  \includegraphics[width=\textwidth]{L#1-oblique.png}\par
  {\small\color{ink}Whole design, oblique view, raking light, thickness direction exaggerated \LadderExaggeration$\times$; colour is height above the mid-plane (scale on the summary page).}\par
  \includegraphics[width=\textwidth]{L#1-blade.png}\par
  {\small\color{ink}Blade close-up (grinds, serrations, tip), same camera rotation, light and colour.}\par
  \includegraphics[width=\textwidth]{L#1-displacement.png}\par
  {\small\color{ink}Top view; colour is displacement from L0 on the shared scale (summary page). L0 is zero everywhere.}\par}
% Generated by smoothing_ladder_figure; do not edit.
\newcommand{\LadderLabel}{\texttt{re\_w406t3\_r01\_p10\_c16}}
\newcommand{\LadderExported}{80.299390}
\newcommand{\LadderExaggeration}{2.5}
\newcommand{\LadderStockHalf}{2.61}
\newcommand{\LadderDispMax}{0.633}

\begin{document}
{\sffamily\bfseries\color{ink}\LARGE Smoothing ladder of the best evolved design}\par
{\sffamily\large \LadderLabel: $W_{\mathrm{ULTRALIGHT\_v1}}=\LadderExported$ as exported, ten smoothing levels}\par
{\small Research draft; author approval pending. Generated by \code{make -C docs/paper smoothing-ladder}; receipt \code{docs/paper/ladder/ladder.json}.}\par
\medskip
\textbf{What this shows.} The exported triangle mesh of the design is smoothed at ten strengths, from none (L0) to the strongest setting the smoother completes on its knife edges without reversing a face (L9). Each level is re-scored with the canonical \code{W_ULTRALIGHT_v1} and rendered from one fixed oblique camera under one low raking light, with the thickness direction exaggerated and faces coloured by height, so that grinds, serrations, channels and faceting stand out. Every camera, light, colour range and exaggeration is shared by all levels; image differences are mesh differences.

\textbf{Smoothing.} Taubin $\lambda/\mu$ filtering ($\lambda=0.5$, $\mu=-0.5$), keeping only the component of each umbrella vector along the area-weighted vertex normal, total displacement capped at 2\,mm. Vertices on edges whose two face normals differ by more than the pin angle do not move. L1--L5 pin at $30^\circ$ (L4, 20 iterations, is the R1 setting of paper variant B); L6--L9 pin at $60^\circ$, the weakest pin angle that completes on this mesh.

\textbf{Re-scoring.} Three of the twenty objectives read the exported 3D geometry and are re-measured at every level: 3D mass (enclosed volume times $4.43\,\mathrm{g\,cm^{-3}}$), the solid tip (scikit-fem P2 submodel on the smoothed skin, cross-checked by DOLFINx~0.9.0 on the same mesh) and the CAD/B-rep objective. The B-rep objective is withheld (scored 0) on every smoothed level, because a smoothed triangle mesh has no validated STEP solid; that costs 5 points at once, so the table also gives $F$ with the B-rep objective kept as exported, which isolates the effect of smoothing. The other seventeen objectives are computed from the genome and carry over unchanged: in particular, edge sharpness and point profile are \emph{not} re-measured on the smoothed mesh. Neither score is a physical test.

\begin{center}
\small
\setlength{\tabcolsep}{4pt}
\begin{tabular}{@{}lrrrrrrrrr@{}}
\toprule
Level & Pins & Iter. & Max disp. & RMS disp. & Vol. & Mass & Tip SF & $F$ & $F$ (B-rep kept) \\
 & & & (mm) & (mm) & (\%) & (g) & (DOLFINx) & & \\
\midrule
% Generated by `jarambit smoothing-ladder`; do not edit.
L0 & -- & 0 & -- & -- & -- & 83.5536 & 0.6886 & 80.623796 & 80.623796 \\
L1 & 30$^\circ$ & 2 & 0.087 & 0.0058 & -0.014 & 83.5420 & 0.6964 & 75.650199 & 80.650199 \\
L2 & 30$^\circ$ & 5 & 0.144 & 0.0096 & -0.031 & 83.5277 & 0.7036 & 75.675500 & 80.675500 \\
L3 & 30$^\circ$ & 10 & 0.177 & 0.0127 & -0.054 & 83.5083 & 0.7040 & 75.681721 & 80.681721 \\
L4 & 30$^\circ$ & 20 & 0.218 & 0.0161 & -0.088 & 83.4803 & 0.7026 & 75.684768 & 80.684768 \\
L5 & 30$^\circ$ & 60 & 0.297 & 0.0231 & -0.149 & 83.4291 & 0.6203 & 75.451120 & 80.451120 \\
L6 & 60$^\circ$ & 20 & 0.503 & 0.0366 & 0.006 & 83.5585 & 0.6957 & 75.643828 & 80.643828 \\

\bottomrule
\end{tabular}
\end{center}

\begin{center}
\begin{tikzpicture}
\begin{axis}[width=0.9\textwidth, height=5.2cm, xlabel={smoothing level}, ylabel={$W_{\mathrm{ULTRALIGHT\_v1}}$}, xtick={0,...,9}, xticklabels={L0,L1,L2,L3,L4,L5,L6,L7,L8,L9},
  legend pos=south west, legend style={font=\small}, grid=major, grid style={gray!25}, tick label style={font=\small}, label style={font=\small}]
\addplot[mark=*, color=teal, thick] table[x=level, y=score_with_brep] {ladder/ladder-scores.dat};
\addlegendentry{$F$, B-rep objective kept as exported}
\addplot[mark=square*, color=rust, thick] table[x=level, y=score] {ladder/ladder-scores.dat};
\addlegendentry{$F$, B-rep objective withheld on smoothed levels}
\end{axis}
\end{tikzpicture}
\end{center}

\textbf{Colour scales.} Height above the mid-plane, 0 to the stock half-thickness \LadderStockHalf\,mm (thicker parts such as the finger ring saturate):\par
\includegraphics[width=\textwidth]{height-colorbar.png}\par
{\small 0\,mm \hfill \LadderStockHalf\,mm}\par
Displacement from L0, shared by all levels, 0 to \LadderDispMax\,mm:\par
\includegraphics[width=\textwidth]{displacement-colorbar.png}\par
{\small 0\,mm \hfill \LadderDispMax\,mm}\par

% Generated by `jarambit smoothing-ladder`; do not edit.
\LadderLevel{0}{0}{0}{80.299390}{80.299390}{0.000}{0.0000}{0.6823}{83.2231}
\LadderLevel{1}{30}{2}{75.309944}{80.309944}{0.101}{0.0056}{0.6848}{83.2113}
\LadderLevel{2}{30}{5}{75.316252}{80.316252}{0.166}{0.0093}{0.6856}{83.1962}
\LadderLevel{3}{30}{10}{75.327832}{80.327832}{0.194}{0.0123}{0.6877}{83.1758}
\LadderLevel{4}{30}{20}{75.343635}{80.343635}{0.214}{0.0156}{0.6905}{83.1472}
\LadderLevel{5}{30}{60}{75.185235}{80.185235}{0.252}{0.0222}{0.6333}{83.0962}
\LadderLevel{6}{60}{20}{75.301958}{80.301958}{0.423}{0.0361}{0.6842}{83.2352}
\LadderLevel{7}{60}{60}{75.290069}{80.290069}{0.495}{0.0502}{0.6750}{83.1745}
\LadderLevel{8}{60}{120}{75.097738}{80.097738}{0.613}{0.0636}{0.5997}{83.0448}
\LadderLevel{9}{60}{200}{75.062571}{80.062571}{0.703}{0.0766}{0.5747}{82.8909}

\end{document}
